% STATUS: reviewed
% Generated from the live database on 13 November 2026 by the documentation's own script
% (the list of acts with their counts) and a table of subject lines that shorten the official
% titles. To refresh it, export the list again and run the script; do not edit the tables by hand.
\chapter{Teisės aktų katalogas}\label{app:acts}

Šiame priede išvardyti visi teisės aktai, kurie 2026~m. lapkričio 13~d. buvo sistemos kataloge: 78 aktai, iš jų 68 įjungti.
Aktas, pažymėtas \chip[muted]{išjungtas}, yra kataloge, tačiau nenaudojamas: jis neatsisiunčiamas,
iš jo neišskiriamos pareigos ir jis nesiūlomas ataskaitoms (kodėl kiekvienas toks aktas pateikiamas
išjungtas, paaiškinta \cref{ch:corpus}). Lentelės suskirstytos pagal šaltinių rūšis. Kiekvienam
įjungtam aktui jos nurodo, ką apima dabartinė jo redakcija: kiek straipsnių rado struktūros
atpažinimas (\enquote{--}, jei tekste straipsnių nėra, pvz., gairėse ar licencijavimo puslapyje),
kiek jame \emph{nuostatų} --- kiekviena nuostata, iki pat punkto, saugoma atskira
eilute, --- kiek pareigų registre yra iš jo išskirtų ir neatmestų pareigų
(\cref{ch:obligations}) ir kiek iš jo nuskaityta sankcijų pakopų (\cref{ch:obligations}). Iš viso
68 įjungtuose aktuose yra 26\,983 nuostatos, 5\,950 pareigų ir 262 sankcijų pakopos.

Trumpasis pavadinimas pirmajame stulpelyje --- tai pavadinimas, kuriuo aktą vadina pati sistema
(ataskaitoje, nuorodoje, pokalbyje), todėl Lietuvos teisės aktų trumpieji pavadinimai yra
lietuviški. Dalykas --- sutrumpintas oficialus akto pavadinimas.

\begin{xltabular}{\textwidth}{@{}L{27mm}YL{21mm}rrrr@{}}
\caption{ES reglamentai, atsisiunčiami iš CELLAR lietuvių kalba: naujausia galiojanti konsoliduota redakcija, o jei jos nėra --- Oficialiajame leidinyje paskelbtas tekstas}\label{tab:acts:eu-reg}\\
\toprule
\hd{Aktas} & \hd{Dalykas} & \hd{CELEX} & \hd{Str.} & \hd{Nuost.} & \hd{Pareig.} & \hd{Pak.} \\
\midrule
\endfirsthead
\toprule
\hd{Aktas} & \hd{Dalykas} & \hd{CELEX} & \hd{Str.} & \hd{Nuost.} & \hd{Pareig.} & \hd{Pak.} \\
\midrule
\endhead
\bottomrule
\endlastfoot
AI Act & Suderintos dirbtinio intelekto taisyklės & \code{32024R1689} & 119 & 1\,513 & 313 & 14 \\
AMLAR & Kovos su pinigų plovimu ir teroristų finansavimu institucijos įsteigimas & \code{32024R1620} & 109 & 977 & 12 & 10 \\
AMLR & Finansų sistemos naudojimo pinigų plovimo ar teroristų finansavimo tikslais prevencija & \code{32024R1624} & 90 & 1\,257 & 271 & 0 \\
CRA & Horizontalieji kibernetinio saugumo reikalavimai produktams su skaitmeniniais elementais (Kibernetinio atsparumo aktas) & \code{32024R2847} & 71 & 757 & 214 & 3 \\
DORA & Finansų sektoriaus skaitmeninis operacinis atsparumas & \code{32022R2554} & 64 & 927 & 204 & 3 \\
DSA & Bendroji skaitmeninių paslaugų rinka (Skaitmeninių paslaugų aktas) & \code{32022R2065} & 93 & 985 & 157 & 7 \\
ECSPR & Europos verslo sutelktinio finansavimo paslaugų teikėjai & \code{32020R1503} & 51 & 751 & 143 & 2 \\
eIDAS 2 & Elektroninė atpažintis ir elektroninių operacijų patikimumo užtikrinimo paslaugos, su pakeitimais, padarytais Reglamentu (ES) 2024/1183 & \code{32014R0910} & 82 & 821 & 158 & 2 \\
GDPR & Fizinių asmenų apsauga tvarkant asmens duomenis & \code{32016R0679} & 99 & 892 & 162 & 7 \\
MAR & Piktnaudžiavimas rinka & \code{32014R0596} & 41 & 596 & 79 & 17 \\
MiCA & Kriptoturto rinkos & \code{32023R1114} & 150 & 2\,296 & 418 & 17 \\
TFR & Informacija, teikiama kartu su lėšų ir tam tikro kriptoturto pervedimais & \code{32023R1113} & 40 & 374 & 76 & 1 \\
CRR \chip[muted]{išjungtas} & Prudenciniai reikalavimai kredito įstaigoms ir investicinėms įmonėms & \code{32013R0575} & -- & -- & -- & -- \\
\end{xltabular}

\begin{xltabular}{\textwidth}{@{}L{27mm}YL{21mm}rrrr@{}}
\caption{Techniniai reguliavimo ir įgyvendinimo standartai pagal MiCA ir DORA}\label{tab:acts:eu-l2}\\
\toprule
\hd{Aktas} & \hd{Dalykas} & \hd{CELEX} & \hd{Str.} & \hd{Nuost.} & \hd{Pareig.} & \hd{Pak.} \\
\midrule
\endfirsthead
\toprule
\hd{Aktas} & \hd{Dalykas} & \hd{CELEX} & \hd{Str.} & \hd{Nuost.} & \hd{Pareig.} & \hd{Pak.} \\
\midrule
\endhead
\bottomrule
\endlastfoot
DORA RTS 2024/1772 & Su IRT susijusių incidentų ir kibernetinių grėsmių klasifikavimas; reikšmingumo ribos & \code{32024R1772} & 13 & 119 & 11 & 0 \\
DORA RTS 2024/1773 & IRT paslaugų, kuriomis remiamos ypatingos svarbos ar svarbios funkcijos, politikos turinys & \code{32024R1773} & 11 & 119 & 36 & 0 \\
DORA RTS 2024/1774 & IRT rizikos valdymo priemonės, metodai, procesai ir politika; supaprastinta sistema & \code{32024R1774} & 42 & 605 & 196 & 0 \\
DORA RTS 2025/301 & Pranešimų ir ataskaitų apie didelius su IRT susijusius incidentus turinys ir terminai & \code{32025R0301} & 7 & 76 & 10 & 0 \\
DORA RTS 2025/532 & Ką nustatyti ir įvertinti, kai IRT paslaugos, kuriomis remiamos ypatingos svarbos ar svarbios funkcijos, perduodamos subrangovams & \code{32025R0532} & 7 & 76 & 12 & 0 \\
MiCA RTS 2025/1140 & Visų kriptoturto paslaugų, veiklos, pavedimų ir sandorių įrašai, kuriuos privaloma saugoti & \code{32025R1140} & 18 & 142 & 64 & 0 \\
MiCA RTS 2025/1142 & Interesų konfliktų politika ir procedūros bei jų atskleidimas & \code{32025R1142} & 9 & 125 & 45 & 0 \\
MiCA RTS 2025/294 & Kriptoturto paslaugų teikėjų skundų nagrinėjimas & \code{32025R0294} & 9 & 64 & 33 & 0 \\
MiCA RTS 2025/299 & Kriptoturto paslaugų teikimo tęstinumas ir reguliarumas & \code{32025R0299} & 7 & 86 & 16 & 0 \\
MiCA RTS 2025/416 & Prekybos platformos pavedimų knygos įrašų turinys ir formatas & \code{32025R0416} & 16 & 80 & 53 & 0 \\
MiCA RTS 2025/417 & Kaip prekybos platforma pateikia skaidrumo duomenis & \code{32025R0417} & 6 & 41 & 15 & 0 \\
MiCA RTS 2025/422 & Tvarumo rodikliai: informacija apie neigiamą poveikį klimatui ir aplinkai & \code{32025R0422} & 7 & 77 & 26 & 0 \\
MiCA RTS 2025/885 & Priemonės, sistemos ir procedūros, skirtos piktnaudžiavimui rinka užkirsti kelią, jį nustatyti ir apie jį pranešti & \code{32025R0885} & 9 & 96 & 24 & 0 \\
DORA ITS 2024/2956 \chip[muted]{išjungtas} & Informacijos registro standartiniai šablonai & \code{32024R2956} & -- & -- & -- & -- \\
DORA ITS 2025/302 \chip[muted]{išjungtas} & Pranešimo apie didelį su IRT susijusį incidentą formos, šablonai ir procedūros & \code{32025R0302} & -- & -- & -- & -- \\
DORA RTS 2025/1190 \chip[muted]{išjungtas} & Grėsmėmis grindžiamas įsiskverbimo testavimas & \code{32025R1190} & -- & -- & -- & -- \\
MiCA ITS 2025/304 \chip[muted]{išjungtas} & Finansų subjekto pranešimo apie ketinimą teikti kriptoturto paslaugas formos, šablonai ir procedūros & \code{32025R0304} & -- & -- & -- & -- \\
MiCA ITS 2025/306 \chip[muted]{išjungtas} & Paraiškos gauti veiklos leidimą formos, šablonai ir procedūros & \code{32025R0306} & -- & -- & -- & -- \\
MiCA RTS 2025/303 \chip[muted]{išjungtas} & Informacija finansų subjekto pranešime apie ketinimą teikti kriptoturto paslaugas & \code{32025R0303} & -- & -- & -- & -- \\
MiCA RTS 2025/305 \chip[muted]{išjungtas} & Informacija paraiškoje gauti kriptoturto paslaugų teikėjo veiklos leidimą & \code{32025R0305} & -- & -- & -- & -- \\
\end{xltabular}

\begin{xltabular}{\textwidth}{@{}L{27mm}YL{21mm}rrrr@{}}
\caption{ES direktyvos}\label{tab:acts:eu-dir}\\
\toprule
\hd{Aktas} & \hd{Dalykas} & \hd{CELEX} & \hd{Str.} & \hd{Nuost.} & \hd{Pareig.} & \hd{Pak.} \\
\midrule
\endfirsthead
\toprule
\hd{Aktas} & \hd{Dalykas} & \hd{CELEX} & \hd{Str.} & \hd{Nuost.} & \hd{Pareig.} & \hd{Pak.} \\
\midrule
\endhead
\bottomrule
\endlastfoot
AMLD6 & Valstybių narių kuriami mechanizmai pinigų plovimo ar teroristų finansavimo prevencijai & \code{32024L1640} & 80 & 876 & 11 & 5 \\
CCD2 & Vartojimo kredito sutartys & \code{32023L2225} & 50 & 566 & 111 & 1 \\
EMD2 & Elektroninių pinigų įstaigų veiklos pradėjimas, vykdymas ir prudencinė priežiūra & \code{32009L0110} & 24 & 122 & 42 & 0 \\
NIS2 & Priemonės aukštam bendram kibernetinio saugumo lygiui visoje Sąjungoje užtikrinti & \code{32022L2555} & 46 & 689 & 43 & 3 \\
PSD2 & Mokėjimo paslaugos vidaus rinkoje & \code{32015L2366} & 118 & 872 & 226 & 0 \\
Whistleblowing Directive & Asmenų, pranešančių apie Sąjungos teisės pažeidimus, apsauga & \code{32019L1937} & 29 & 389 & 22 & 1 \\
CRD \chip[muted]{išjungtas} & Galimybė verstis kredito įstaigų veikla ir kredito įstaigų prudencinė priežiūra & \code{32013L0036} & -- & -- & -- & -- \\
\end{xltabular}

\begin{xltabular}{\textwidth}{@{}L{27mm}Yrrrr@{}}
\caption{Europos bankininkystės institucijos gairės, įkeltos kaip failai}\label{tab:acts:eba}\\
\toprule
\hd{Aktas} & \hd{Dalykas} & \hd{Str.} & \hd{Nuost.} & \hd{Pareig.} & \hd{Pak.} \\
\midrule
\endfirsthead
\toprule
\hd{Aktas} & \hd{Dalykas} & \hd{Str.} & \hd{Nuost.} & \hd{Pareig.} & \hd{Pak.} \\
\midrule
\endhead
\bottomrule
\endlastfoot
EBA/GL/2021/02 & Pinigų plovimo ir teroristų finansavimo rizikos veiksnių gairės: deramas klientų tikrinimas ir veiksniai, į kuriuos atsižvelgiama vertinant pinigų plovimo ir teroristų finansavimo riziką (konsoliduota redakcija) & -- & 411 & 352 & 0 \\
EBA/GL/2024/01 & Gairės, kuriomis iš dalies keičiamos EBA/GL/2021/02 ir kriptoturto paslaugų teikėjai įtraukiami į rizikos veiksnių gaires & -- & 68 & 45 & 0 \\
\end{xltabular}

\begin{xltabular}{\textwidth}{@{}L{27mm}Yrrrr@{}}
\caption{Lietuvos įstatymai, atsisiunčiami iš Teisės aktų registro (e-TAR) tą dieną galiojančia redakcija}\label{tab:acts:lt-law}\\
\toprule
\hd{Aktas} & \hd{Dalykas} & \hd{Str.} & \hd{Nuost.} & \hd{Pareig.} & \hd{Pak.} \\
\midrule
\endfirsthead
\toprule
\hd{Aktas} & \hd{Dalykas} & \hd{Str.} & \hd{Nuost.} & \hd{Pareig.} & \hd{Pak.} \\
\midrule
\endhead
\bottomrule
\endlastfoot
ADTAĮ & Asmens duomenų teisinės apsaugos įstatymas & 40 & 283 & 27 & 2 \\
BĮ & Bankų įstatymas & 111 & 1\,096 & 294 & 24 \\
EPEPĮĮ & Elektroninių pinigų ir elektroninių pinigų įstaigų įstatymas & 48 & 450 & 127 & 14 \\
FĮĮ & Finansų įstaigų įstatymas & 55 & 517 & 116 & 0 \\
KRĮ & Kriptoturto rinkų įstatymas & 22 & 206 & 50 & 23 \\
KSĮ & Kibernetinio saugumo įstatymas & 38 & 547 & 67 & 4 \\
KUĮ & Kredito unijų įstatymas & 87 & 793 & 309 & 17 \\
LBĮ & Lietuvos banko įstatymas & 83 & 595 & 29 & 13 \\
MĮ & Mokėjimų įstatymas & 98 & 800 & 274 & 19 \\
MĮĮ & Mokėjimo įstaigų įstatymas & 42 & 420 & 119 & 14 \\
NTSKĮ & Su nekilnojamuoju turtu susijusio kredito įstatymas & 62 & 617 & 203 & 8 \\
PAĮ & Pranešėjų apsaugos įstatymas & 17 & 188 & 20 & 0 \\
PPTFPĮ & Pinigų plovimo ir teroristų finansavimo prevencijos įstatymas & 62 & 836 & 207 & 24 \\
VKĮ & Vartojimo kredito įstatymas & 59 & 523 & 143 & 7 \\
\end{xltabular}

\begin{xltabular}{\textwidth}{@{}L{27mm}Yrrrr@{}}
\caption{Lietuvos banko ir Finansinių nusikaltimų tyrimo tarnybos teisės aktai}\label{tab:acts:lt-sup}\\
\toprule
\hd{Aktas} & \hd{Dalykas} & \hd{Str.} & \hd{Nuost.} & \hd{Pareig.} & \hd{Pak.} \\
\midrule
\endfirsthead
\toprule
\hd{Aktas} & \hd{Dalykas} & \hd{Str.} & \hd{Nuost.} & \hd{Pareig.} & \hd{Pak.} \\
\midrule
\endhead
\bottomrule
\endlastfoot
FNTT V-314 & Finansinių nusikaltimų tyrimo tarnyba: techniniai reikalavimai kliento tapatybei nustatyti nuotoliniu būdu tiesioginės vaizdo transliacijos priemonėmis (2016~m. lapkričio 30~d. įsakymas) & -- & 66 & 26 & 0 \\
LB 03-126 & Baudų apskaičiavimo tvarkos aprašas (2018~m. liepos 10~d.) & -- & 65 & 0 & 0 \\
LB 03-17 & Pinigų plovimo ir (ar) teroristų finansavimo prevencijos instrukcija finansų rinkos dalyviams (2015~m. vasario 12~d. valdybos nutarimas) & -- & 168 & 87 & 0 \\
LB 03-175 & Teisės akto pažeidimo pripažinimo mažareikšmiu tvarka (2020~m. lapkričio 26~d.) & -- & 20 & 1 & 0 \\
LB 03-98 & Tarptautinių sankcijų įgyvendinimo instrukcija finansų rinkos dalyviams (2023~m. gegužės 30~d.) & -- & 96 & 43 & 0 \\
LB atitikties funkcijos gairės (2022) & Lietuvos banko atitikties funkcijos gairės (2022) & -- & 200 & 82 & 0 \\
LB 03-174 \chip[muted]{išjungtas} & IRT ir saugumo rizikos valdymo reikalavimai (2020~m. lapkričio 26~d.); panaikinti & -- & -- & -- & -- \\
\end{xltabular}

\begin{xltabular}{\textwidth}{@{}L{27mm}Yrrrr@{}}
\caption{Lietuvos banko licencijavimo puslapiai, po vieną kiekvienai įstaigų rūšiai}\label{tab:acts:lt-lic}\\
\toprule
\hd{Aktas} & \hd{Dalykas} & \hd{Str.} & \hd{Nuost.} & \hd{Pareig.} & \hd{Pak.} \\
\midrule
\endfirsthead
\toprule
\hd{Aktas} & \hd{Dalykas} & \hd{Str.} & \hd{Nuost.} & \hd{Pareig.} & \hd{Pak.} \\
\midrule
\endhead
\bottomrule
\endlastfoot
Banko steigimas ir licencijavimas & Bankų licencijavimas & -- & 61 & 20 & 0 \\
Elektroninių pinigų įstaigos steigimas ir licencijavimas & Elektroninių pinigų įstaigų licencijavimas & -- & 41 & 10 & 0 \\
Elektroninių pinigų žetonų emitentų licencijavimas & Elektroninių pinigų žetonų emitentų licencijavimas & -- & 32 & 4 & 0 \\
Finansų maklerio įmonės steigimas ir licencijavimas & Finansų maklerio įmonių licencijavimas & -- & 75 & 31 & 0 \\
Kito kriptoturto emitentų licencijavimas & Kito kriptoturto emitentų licencijavimas & -- & 29 & 12 & 0 \\
Kredito unijų steigimas ir licencijavimas & Kredito unijų licencijavimas & -- & 94 & 3 & 0 \\
Kriptoturto paslaugų teikėjų licencijavimas & Kriptoturto paslaugų teikėjų licencijavimas & -- & 57 & 2 & 0 \\
Mokėjimo įstaigų steigimas ir licencijavimas & Mokėjimo įstaigų licencijavimas & -- & 57 & 9 & 0 \\
NT kredito davėjai & Kredito davėjai pagal su nekilnojamuoju turtu susijusio kredito įstatymą & -- & 23 & 2 & 0 \\
Su turtu susietų žetonų emitentų licencijavimas & Su turtu susietų žetonų emitentų licencijavimas & -- & 50 & 4 & 0 \\
Sutelktinio finansavimo paslaugų teikėjai & Sutelktinio finansavimo paslaugų teikėjai & -- & 39 & 13 & 0 \\
Tarpusavio skolinimo platformos operatoriai & Tarpusavio skolinimo platformų operatoriai & -- & 28 & 5 & 0 \\
Valiutos keityklos operatoriai & Valiutos keityklų operatoriai & -- & 24 & 5 & 0 \\
Vartojimo kredito davėjai & Vartojimo kredito davėjai & -- & 26 & 6 & 0 \\
Vartojimo kredito tarpininkai & Vartojimo kredito tarpininkai & -- & 16 & 0 & 0 \\
\end{xltabular}
